% ============================================================
%  Part IV — Prepositions
% ============================================================
\gpart{IV}{Prepositions}

% ------------------------------------------------------------
\gsec{Prepositions by case}

A preposition fixes the case of the noun that follows it. Learn the preposition
and its case as one item; there is no way to work it out from the meaning.

\begin{gtbl}{colspec={X[0.85,l] X[2.15,l]}}
  Case & Prepositions \\
  Akkusativ & durch, für, gegen, ohne, um, bis, entlang, wider \\
  Dativ & aus, außer, bei, gegenüber, mit, nach, seit, von, zu, ab \\
  Wechsel & an, auf, hinter, in, neben, über, unter, vor, zwischen \\
  Genitiv & (an)statt, trotz, während, wegen, außerhalb, innerhalb, oberhalb, unterhalb, diesseits, jenseits, aufgrund \\
\end{gtbl}

\begin{gnote}
Two mnemonics worth the ten seconds. Accusative: \textbf{DOGFUB} — durch, ohne,
gegen, für, um, bis. Dative: chant it in order — \textit{aus, außer, bei,
gegenüber, mit, nach, seit, von, zu}.
\end{gnote}

\tcap{Contractions}

\begin{flattbl}{colspec={X[1,c] X[1,c] X[1,c] X[1,c]}}
  an + dem = am    & an + das = ans   & in + dem = im    & in + das = ins \\
  zu + dem = zum   & zu + der = zur   & bei + dem = beim & von + dem = vom \\
  auf + das = aufs & für + das = fürs & um + das = ums   & durch + das = durchs \\
\end{flattbl}

\begin{gnote}
The contractions are compulsory in speech unless you are pointing at a specific
thing. \textit{Ich gehe ins Kino} is normal; \textit{Ich gehe in das Kino dort
drüben} keeps the article separate for emphasis.
\end{gnote}

% ------------------------------------------------------------
\gsec{Accusative prepositions}

\begin{gtbl}{colspec={X[0.75,l] X[1.1,l] X[1.55,l]}}
  Preposition & Meaning & Example \\
  durch   & through, by means of        & \textbf{durch den} Park \\
  für     & for                         & \textbf{für meinen} Bruder \\
  gegen   & against; around \eng{time}  & \textbf{gegen die} Wand \\
  ohne    & without                     & \textbf{ohne einen} Plan \\
  um      & around; at \eng{clock time} & \textbf{um den} Tisch \\
  bis     & until, as far as            & \textbf{bis nächsten} Freitag \\
  entlang & along \eng{follows the noun} & die Straße \textbf{entlang} \\
  wider   & contrary to \eng{formal}    & \textbf{wider besseres} Wissen \\
\end{gtbl}

\begin{gnote}
\textit{ohne} normally drops the article entirely: \textit{ohne Auto},
\textit{ohne Geld}. \textit{bis} usually appears without an article too, or
paired with another preposition which then sets the case:
\textit{bis \textbf{zum} Bahnhof} is dative.
\end{gnote}

\tcap{In time expressions}

\begin{flattbl}{colspec={X[0.8,l] X[2.2,l]}}
  um    & \textit{um acht Uhr} \eng{at eight} \\
  gegen & \textit{gegen acht} \eng{around eight} \\
  bis   & \textit{bis Montag} \eng{until Monday} \\
  für   & \textit{für zwei Wochen} \eng{for a planned span} \\
\end{flattbl}

% ------------------------------------------------------------
\gsec{Dative prepositions}

\begin{gtbl}{colspec={X[0.85,l] X[1.1,l] X[1.45,l]}}
  Preposition & Meaning & Example \\
  aus       & out of; from \eng{origin}      & \textbf{aus der} Schweiz \\
  außer     & except for                     & \textbf{außer meinem} Bruder \\
  bei       & at, near, at the home of       & \textbf{beim} Arzt \\
  gegenüber & opposite \eng{often follows}   & dem Bahnhof \textbf{gegenüber} \\
  mit       & with; by \eng{transport}       & \textbf{mit dem} Bus \\
  nach      & after; to \eng{places}         & \textbf{nach der} Arbeit \\
  seit      & since, for \eng{time so far}   & \textbf{seit einem} Jahr \\
  von       & from, of, by                   & \textbf{vom} Bahnhof \\
  zu        & to \eng{people, places}        & \textbf{zum} Supermarkt \\
  ab        & from \eng{time onwards}        & \textbf{ab nächster} Woche \\
\end{gtbl}

\begin{gnote}
\textbf{nach vs zu.} Use \textit{nach} with countries, cities and continents
that take no article, and with \textit{Hause}. Use \textit{zu} with people,
businesses and named buildings: \textit{zum Arzt}, \textit{zur Bank},
\textit{nach Berlin}, \textit{nach Hause} \eng{homewards}, but \textit{zu Hause}
\eng{at home}.
\end{gnote}

\begin{gnote}
\textit{seit} takes the present tense in German where English uses the perfect:
\textit{Ich wohne seit drei Jahren hier} — \emph{I have been living here for
three years}.
\end{gnote}

% ------------------------------------------------------------
\gsec{Two-way prepositions}

Nine prepositions take \textbf{either} case: accusative for movement into a new
place, dative for position or for movement within one place.

\begin{gtbl}{colspec={X[0.75,l] X[0.95,l] X[1.25,l] X[1.25,l]}}
  Prep. & Meaning & Wohin? \eng{Akk.} & Wo? \eng{Dat.} \\
  an       & on, at \eng{vertical}  & an \textbf{die} Wand     & an \textbf{der} Wand \\
  auf      & on \eng{horizontal}    & auf \textbf{den} Tisch   & auf \textbf{dem} Tisch \\
  hinter   & behind                 & hinter \textbf{das} Haus & hinter \textbf{dem} Haus \\
  in       & in, into               & in \textbf{die} Stadt    & in \textbf{der} Stadt \\
  neben    & next to                & neben \textbf{den} Stuhl & neben \textbf{dem} Stuhl \\
  über     & over, above            & über \textbf{die} Straße & über \textbf{der} Straße \\
  unter    & under, among           & unter \textbf{den} Tisch & unter \textbf{dem} Tisch \\
  vor      & in front of; ago       & vor \textbf{die} Tür     & vor \textbf{der} Tür \\
  zwischen & between                & zwischen \textbf{die} Bäume & zwischen \textbf{den} Bäumen \\
\end{gtbl}

\tcap{The test}

\begin{flattbl}{colspec={X[0.75,l] X[0.85,l] X[1.9,l]}}
  \textbf{Wohin?} & Akkusativ & crossing a boundary into somewhere new \\
  \textbf{Wo?}    & Dativ     & already there, or moving about inside it \\
\end{flattbl}

\begin{gnote}
The verb usually gives it away. \textit{stellen, legen, setzen, hängen}
\eng{to put} take the accusative; their partners \textit{stehen, liegen, sitzen,
hängen} \eng{to be} take the dative. \textit{Ich hänge das Bild an die Wand}
versus \textit{Das Bild hängt an der Wand}.
\end{gnote}

\begin{gnote}
Used non-spatially these prepositions are fixed, not free: \textit{an etwas
denken} and \textit{sich auf etwas freuen} always take the accusative;
\textit{vor einem Jahr} \eng{a year ago} always the dative.
\end{gnote}

% ------------------------------------------------------------
\gsec{Genitive prepositions}

\begin{gtbl}{colspec={X[0.95,l] X[1,l] X[1.3,l]}}
  Preposition & Meaning & Example \\
  während    & during          & \textbf{während des} Films \\
  wegen      & because of      & \textbf{wegen des} Wetters \\
  trotz      & in spite of     & \textbf{trotz der} Kälte \\
  (an)statt  & instead of      & \textbf{statt eines} Briefes \\
  außerhalb  & outside of      & \textbf{außerhalb der} Stadt \\
  innerhalb  & within          & \textbf{innerhalb einer} Woche \\
  oberhalb   & above           & \textbf{oberhalb des} Dorfes \\
  unterhalb  & below           & \textbf{unterhalb des} Gipfels \\
  diesseits  & this side of    & \textbf{diesseits der} Grenze \\
  jenseits   & beyond          & \textbf{jenseits des} Flusses \\
  aufgrund   & on the basis of & \textbf{aufgrund der} Daten \\
  angesichts & in view of      & \textbf{angesichts der} Lage \\
\end{gtbl}

\begin{gnote}
In speech \textit{wegen}, \textit{trotz}, \textit{während} and \textit{statt}
are very often heard with the dative (\textit{wegen dem Wetter}). It is common
but still non-standard, so keep the genitive in writing.
\end{gnote}

\begin{gnote}
\textit{wegen} fuses with pronouns into fixed forms: \textit{meinetwegen},
\textit{deinetwegen}, \textit{seinetwegen}, \textit{unsertwegen}.
\textit{Meinetwegen} on its own means \emph{fine by me}.
\end{gnote}
